\section{Interval configurations, labelled relations, and torsion}
\label{sec:states-filtration}

Let $\mathcal C$ be the set of finite subsets $S\subset\N_{>0}$ whose connected components in the path graph have cardinality $2$ or $3$.  Put
\begin{equation}\label{eq:observations}
W(S)=\sum_{n\in S}\frac1n,
\qquad
\nu(S)=|\pi_0(S)|,
\qquad
\overline W(S)=[W(S)]\in A:=\Q/\Z.
\end{equation}
For $n\ge1$, let
\[
D_n\subseteq\mathcal C^n
\]
be the joint-compatibility domain: $(S_1,\ldots,S_n)\in D_n$ when the supports are pairwise disjoint and their literal union is again an element of $\mathcal C$.  On $D_n$ the partial product is simply union.  Hence
\begin{equation}\label{eq:additive-observations}
W\!\left(\bigcup_iS_i\right)=\sum_iW(S_i),
\quad
\nu\!\left(\bigcup_iS_i\right)=\sum_i\nu(S_i),
\quad
\overline W\!\left(\bigcup_iS_i\right)=\sum_i\overline W(S_i).
\end{equation}
Equivalently, for binary compatible unions the diagram
\begin{equation}\label{eq:observation-square}
\begin{CD}
D_2 @>{\cup}>> \mathcal C\\
@V{(\Omega\pi_1,\Omega\pi_2)}VV @VV{\Omega}V\\
(A\times\N)^2 @>{+}>> A\times\N
\end{CD}
\qquad
\Omega=(\overline W,\nu)
\end{equation}
commutes.

\subsection{The locally posetal 2-category of labelled relations}

For finite sets $X,Y$, a \emph{labelled finite relation} $R:X\rightsquigarrow Y$ is a finite subset
\[
R\subseteq X\times\mathcal C\times Y.
\]
It is \emph{covering} if its projection to $Y$ is surjective.  Thus covering is a property of a labelled relation, not part of the definition of the class of morphisms.

For labelled relations $R:X\rightsquigarrow Y$ and $S:Y\rightsquigarrow Z$, define $S\odot R:X\rightsquigarrow Z$ by
\[
(x,L,z)\in S\odot R
\]
if and only if there are $y\in Y$ and labels $L_1,L_2\in\mathcal C$ such that
\[
(x,L_1,y)\in R,
\qquad
(y,L_2,z)\in S,
\qquad
(L_1,L_2)\in D_2,
\qquad
L=L_1\cup L_2.
\]
Equivalently, $S\odot R$ is obtained from the fibre product of the two edge sets over $Y$, restricted to compatible label pairs, by applying literal union to the labels and then taking the image.  Put
\[
1_X=\{(x,\varnothing,x):x\in X\}.
\]

\begin{lemma}\label{lem:labelled-relation-category}
The operation $\odot$ is associative and $1_X$ is a two-sided identity.  With inclusions of labelled relations as $2$-cells, finite sets and labelled finite relations therefore form a locally posetal $2$-category, denoted $\mathbf{Rel}_{\mathcal C}$.
\end{lemma}

\begin{proof}
The empty configuration is compatible with every label, so $1_X$ is an identity.  Suppose $L_1,L_2,L_3\in\mathcal C$ are pairwise disjoint and their total union belongs to $\mathcal C$.  No component of one $L_i$ can be adjacent to a component of another: such two components each have cardinality at least $2$, so their union would lie in one connected component of cardinality at least $4$.  Hence every partial union of the $L_i$ also belongs to $\mathcal C$.  It follows that either parenthesization of a triple composite selects exactly the same triples of edges, and both assign the literal union $L_1\cup L_2\cup L_3$ as label.  Thus $\odot$ is associative.  Monotonicity under inclusion is immediate.
\end{proof}

The underlying unlabelled relation is
\begin{equation}\label{eq:underlying-relation}
U(R)=\{(x,y):\exists L\in\mathcal C,\ (x,L,y)\in R\}\subseteq X\times Y.
\end{equation}
For $R\ne\varnothing$ put
\[
\omega(R)=\max\{W(L):(x,L,y)\in R\},
\]
and set $\omega(\varnothing)=0$.

For every composable pair there is a canonical inclusion
\begin{equation}\label{eq:oplax-inclusion}
U(S\odot R)\subseteq U(S)\circ U(R).
\end{equation}
Thus $U$ preserves identities strictly and carries composition to the comparison $2$-cell
\begin{equation}\label{eq:oplax-2cell}
U(S\odot R)
\xRightarrow{\;\vartheta_{S,R}\;}
U(S)\circ U(R),
\end{equation}
where a $2$-cell in $\mathbf{Rel}$ is inclusion of relations.  Associativity of $\odot$ makes these comparisons coherent automatically; in other words,
\[
U:\mathbf{Rel}_{\mathcal C}\longrightarrow\mathbf{Rel}
\]
is a normal oplax $2$-functor.

For a finite composable string $R_1,\ldots,R_n$, let $P_n$ be the fibre product of the edge sets over the intermediate sets.  The labels define a map $\ell:P_n\to\mathcal C^n$.  The obstruction to strictness is exactly whether $\ell$ lands in $D_n$:
\begin{equation}\label{eq:compatibility-factorization}
\begin{CD}
P_n @>{\ell}>> \mathcal C^n\\
@V{\widetilde\ell}VV @AA{\iota}A\\
D_n @= D_n
\end{CD}
\end{equation}
Whenever such a factorization exists, every relationally composable witness is jointly compatible, and hence
\begin{equation}\label{eq:rel-strictification}
U(R_n\odot\cdots\odot R_1)
=
U(R_n)\circ\cdots\circ U(R_1).
\end{equation}
Moreover, for every composable string,
\begin{equation}\label{eq:mass-subadditive}
\omega(R_n\odot\cdots\odot R_1)
\le\sum_{i=1}^n\omega(R_i),
\end{equation}
because reciprocal value is additive on jointly compatible unions.

\subsection{Finite torsion in \texorpdfstring{$\Q/\Z$}{Q/Z}}

For $n\ge1$ put
\[
H_n=\{x\in A:nx=0\}.
\]
For a subgroup $H\le A$, let
\[
q_H:A\longrightarrow A/H
\]
be the quotient map, and for $\alpha\in A/H$ write
\[
\operatorname{Tor}_H(\alpha)=q_H^{-1}(\alpha)
\]
for the corresponding affine torsor fibre.

\begin{lemma}\label{lem:torsion-filtration}
For $m,n\ge1$,
\[
H_n=\langle[1/n]\rangle,
\qquad |H_n|=n,
\]
\[
H_m\le H_n\iff m\mid n,
\qquad
H_m\vee H_n=H_{\operatorname{lcm}(m,n)}.
\]
Consequently the nonzero join-irreducible stages are exactly $H_{p^e}$, with $p$ prime and $e\ge1$.
\end{lemma}

\begin{proof}
If $[q]\in H_n$, then $nq\in\Z$, so $[q]=[j/n]$ for some $j\in\Z$; hence $H_n=\langle[1/n]\rangle$ and has order $n$.  The inclusion criterion follows by testing $[1/m]$.  Put $L=\operatorname{lcm}(m,n)$.  Since $\gcd(L/m,L/n)=1$, there exist $a,b\in\Z$ with
\[
aL/m+bL/n=1.
\]
Dividing by $L$ gives $[1/L]\in H_m+H_n$, whence $H_m\vee H_n=H_L$.  The join-irreducibles of the divisibility lattice are precisely the prime powers.
\end{proof}

Accordingly we index the nontrivial steps by prime powers $Q=p^e$.  Let $F_Q$ be the join of all prime-power stages of rank at most $Q$, and $F_{<Q}$ the join of those of smaller rank.  The simple quotient
\begin{equation}\label{eq:simple-quotient}
S_Q=F_Q/F_{<Q}
\end{equation}
has order $p$: the $p$-adic valuation of the lower least common multiple is $e-1$, and adjoining $p^e$ raises it by one.  Thus $S_Q$ is canonically a one-dimensional $\F_p$-object.

For $X\ge1$ define the endpoint stage
\begin{equation}\label{eq:endpoint-stage}
E_X=H_{\operatorname{lcm}(1,\ldots,\lfloor X\rfloor)}.
\end{equation}
Every class $[a/b]\in A$ belongs to $H_b$, so the filtration $(E_X)$ is cofinal and every fixed class of $A$ becomes zero in $A/E_X$ for all sufficiently large $X$.

Finally, if $Q$ is a prime-power step, then
\begin{equation}\label{eq:endpoint-current-step}
E_{Q-1}=F_{<Q},
\qquad
E_Q=F_Q.
\end{equation}
Indeed a prime-power rank is $<Q$ exactly when it is at most $Q-1$.  Thus the local simple quotient \eqref{eq:simple-quotient} is literally the associated quotient of the global endpoint filtration at $Q$.
